% !TEX root = ../main.tex

\section{Softvérová časť}

Softvérová časť projektu zabezpečuje prepojenie medzi hardvérovým modelom, 
lokálnym webovým rozhraním, databázou a cloudovou platformou ThingsBoard. 
Hlavnou úlohou softvéru je prijímať dáta z mikrokontroléra ESP32, 
spracovať aktuálny stav systému, zobraziť namerané hodnoty v reálnom čase, 
ukladať historické merania a umožniť lokálne aj vzdialené ovládanie systému.

Základný tok dát je navrhnutý tak, že ESP32 periodicky odosiela aktuálne hodnoty na 
Flask server. Server prijaté dáta spracuje, odošle ich do lokálneho webového dashboardu 
pomocou WebSocket komunikácie, pri aktívnom meraní ich uloží do dočasného buffra 
a zároveň publikuje telemetriu do platformy ThingsBoard. Po zastavení merania 
sa celá meracia sekvencia uloží do MySQL databázy a do záložného JSON súboru.

\begin{figure}[H]
    \centering
    \safeimage{obrazky/tok_dat.png}{0.90\textwidth}{Diagram toku dát softvérovej časti systému}
    \label{fig:tok_dat_softver}
\end{figure}

\subsection{Architektúra softvérového riešenia}

Softvérové riešenie je sústredené okolo Flask servera, ktorý prepája ESP32, webové rozhranie, databázu a platformu ThingsBoard. ESP32 zabezpečuje meranie a výkonové riadenie, zatiaľ čo server spracúva prijaté dáta, uchováva aktuálny stav systému a poskytuje rozhrania pre lokálne aj vzdialené ovládanie.

Hlavné softvérové časti sú:

\begin{itemize}
    \item \textbf{ESP32 firmvér} -- odosielanie nameraných hodnôt a prijímanie riadiacich parametrov,
    \item \textbf{Flask server} -- REST API, WebSocket komunikácia, stav systému a ukladanie dát,
    \item \textbf{webové rozhranie} -- realtime grafy, ovládacie prvky a archív meraní,
    \item \textbf{ThingsBoard integrácia} -- MQTT telemetria a vybrané vzdialené príkazy.
\end{itemize}

\subsection{Komunikácia ESP32 so serverom}

Mikrokontrolér ESP32 odosiela namerané údaje na Flask server pomocou HTTP požiadavky 
na endpoint \path{/api/update}. Odosielaná správa obsahuje najmä aktuálnu teplotu kvapaliny, 
PWM hodnotu pumpy a PWM hodnotu Peltierovho článku.

\begin{lstlisting}[language=json, caption={Príklad dát odosielaných z ESP32 na server}]
{
  "temperature": 24.5,
  "pump_pwm": 150,
  "tec_pwm": 120
}
\end{lstlisting}

Server po prijatí dát aktualizuje svoj interný stav a odošle ESP32 spätnú odpoveď. 
Tá obsahuje napríklad informáciu o tom, či je systém pripojený, či beží meranie, 
aký regulačný mód je zvolený, aké sú cieľové PWM hodnoty a aká je nastavená požadovaná teplota.

\begin{lstlisting}[language=json, caption={Príklad odpovede servera pre ESP32}]
{
  "status": "ok",
  "connected": true,
  "running": true,
  "mode": 1,
  "target_pump": 150,
  "target_tec": 120,
  "setpoint": 25.0
}
\end{lstlisting}

Dôležitým prvkom je stav \path{connected}. Ak server vráti hodnotu \path{connected=false}, 
hardvér ostáva v bezpečnom stave a výkonové prvky nemusia byť aktivované. Tým sa zabraňuje 
nechcenému spusteniu Peltierovho článku alebo pumpy.

\subsection{Backend Flask server a REST API}

Flask server tvorí centrálny bod softvérovej časti. Prijíma dáta z ESP32, uchováva 
aktuálny stav systému, spracúva príkazy z webového rozhrania, komunikuje s ThingsBoardom a 
zabezpečuje ukladanie meraní. Aktuálne hodnoty sú uložené v stavovej premennej 
\path{system_state}, ktorá obsahuje napríklad teplotu, PWM hodnoty, setpoint, regulačný mód, 
stav merania a buffer aktuálnej session.

Pre komunikáciu sú vytvorené REST API endpointy uvedené v tabuľke \ref{tab:rest_api_endpointy}.

{\small
\begin{longtable}{|>{\RaggedRight\arraybackslash}p{0.33\textwidth}|>{\centering\arraybackslash}p{0.12\textwidth}|>{\RaggedRight\arraybackslash}p{0.45\textwidth}|}
\hline
\textbf{Endpoint} & \textbf{Metóda} & \textbf{Úloha} \\
\hline
\path{/api/update} & POST & Príjem aktuálnych dát z ESP32. \\
\hline
\path{/api/connect} & POST & Logické pripojenie alebo odpojenie systému. \\
\hline
\path{/api/start} & POST & Spustenie merania alebo regulácie. \\
\hline
\path{/api/stop} & POST & Zastavenie systému a uloženie aktuálnej session. \\
\hline
\path{/api/setpoint} & POST & Nastavenie požadovanej teploty. \\
\hline
\path{/api/mode} & POST & Nastavenie regulačného módu. \\
\hline
\path{/api/pwm_pump} & POST & Nastavenie cieľovej PWM hodnoty pumpy. \\
\hline
\path{/api/pwm_tec} & POST & Nastavenie cieľovej PWM hodnoty Peltierovho článku. \\
\hline
\path{/api/read_db/<id>} & GET & Čítanie historického záznamu z databázy. \\
\hline
\path{/api/read_log/<id>} & GET & Čítanie historického záznamu zo súboru JSON. \\
\hline
\caption{Prehľad hlavných REST API endpointov}
\label{tab:rest_api_endpointy}
\end{longtable}
}

Endpoint \path{/api/update} je hlavný vstupný bod pre dáta z hardvéru. Endpointy \path{/api/start} a \path{/api/stop} riadia začiatok a koniec merania. Endpointy pre setpoint, mód a PWM hodnoty slúžia na zmenu riadiacich parametrov systému.

\subsection{Webové rozhranie a realtime prenos dát}

Lokálne webové rozhranie slúži na sledovanie aktuálneho stavu systému a na zadávanie 
riadiacich príkazov. Frontend je vytvorený pomocou HTML, CSS a JavaScriptu. 
Grafy sú vykresľované pomocou knižnice Plotly.

Realtime prenos dát medzi serverom a webovým rozhraním je realizovaný pomocou WebSocket 
komunikácie cez Flask-SocketIO. Po prijatí nových dát zo zariadenia server odošle udalosť 
\path{new_data}, ktorú frontend použije na aktualizáciu grafov.

Webové rozhranie zobrazuje najmä:

\begin{itemize}
    \item aktuálnu teplotu kvapaliny,
    \item požadovanú teplotu,
    \item PWM hodnotu pumpy,
    \item PWM hodnotu Peltierovho článku,
    \item aktuálny stav systému a regulačný mód.
\end{itemize}

Používateľ môže cez webové rozhranie systém pripojiť, spustiť alebo zastaviť meranie, 
nastaviť požadovanú teplotu, zvoliť regulačný mód a manuálne nastaviť PWM hodnoty pumpy a 
Peltierovho článku. Súčasťou webovej aplikácie je aj archívna stránka, ktorá umožňuje 
zobrazenie historických meraní z databázy alebo zo súboru JSON.

\subsection{Dátová vrstva -- MySQL a JSON}

Dátová vrstva je navrhnutá ako dvojúrovňové ukladanie nameraných údajov. 
Počas aktívneho merania sa dáta ukladajú do dočasného buffra \path{session_buffer}. 
Dátové body sa do buffra pridávajú iba vtedy, keď je systém v stave \path{running=True}. 
Tým sa zabraňuje ukladaniu zbytočných údajov v pohotovostnom režime.

Každý dátový bod obsahuje najmä poradové číslo vzorky, časovú značku, teplotu, 
PWM hodnotu pumpy, PWM hodnotu Peltierovho článku, požadovanú teplotu a aktuálny regulačný mód.

Po zastavení merania sa celá session uloží naraz do dvoch úložísk:

\begin{itemize}
    \item do MySQL databázy \path{ml_poit}, tabuľka \path{graph},
    \item do lokálneho súboru \path{data_log.json}.
\end{itemize}

Databáza slúži ako hlavné úložisko historických meraní. JSON súbor slúži ako 
jednoduchá záloha a umožňuje čítanie dát aj bez databázového servera. 
Po úspešnom uložení sa buffer vyčistí a systém je pripravený na ďalšie meranie.

\subsection{ThingsBoard integrácia}

Pre vzdialený monitoring a manažment systému bola použitá platforma ThingsBoard. 
Flask server komunikuje s ThingsBoardom pomocou MQTT protokolu na porte 
\path{1883}. Po prijatí dát z ESP32 server publikuje telemetriu na topic 
\path{v1/devices/me/telemetry}.

Do ThingsBoardu sa odosielajú hlavné veličiny uvedené v tabuľke \ref{tab:thingsboard_telemetria}.

{\small
\begin{longtable}{|>{\RaggedRight\arraybackslash}p{0.26\textwidth}|>{\RaggedRight\arraybackslash}p{0.48\textwidth}|>{\RaggedRight\arraybackslash}p{0.16\textwidth}|}
\hline
\textbf{Kľúč} & \textbf{Popis} & \textbf{Jednotka / rozsah} \\
\hline
\path{temperature} & Nameraná teplota kvapaliny. & $^\circ$C \\
\hline
\path{setpoint} & Požadovaná teplota. & $^\circ$C \\
\hline
\path{error} & Regulačná odchýlka. & $^\circ$C \\
\hline
\path{pump_pwm} & PWM hodnota pumpy. & 0--255 \\
\hline
\path{tec_pwm} & PWM hodnota Peltierovho článku. & 0--255 \\
\hline
\path{mode} & Aktívny regulačný mód. & 1 / 2 / 3 \\
\hline
\path{running} & Stav regulácie. & bool \\
\hline
\caption{Telemetrické veličiny odosielané do ThingsBoardu}
\label{tab:thingsboard_telemetria}
\end{longtable}
}

Okrem číselných hodnôt sa do ThingsBoardu odosiela aj textový atribút \path{terminal}, 
ktorý zobrazuje stručný diagnostický stav systému. Obsahuje napríklad aktuálnu teplotu, 
regulačnú odchýlku, stav RUN/STOP a aktívny mód.

Dashboard v ThingsBoarde je rozdelený na dva stavy. Stav \textit{Default} slúži na realtime 
sledovanie a ovládanie systému. Obsahuje grafy teploty, PWM pumpy a PWM Peltierovho článku, 
tlačidlá \textit{Connect}, \textit{Start}, \textit{Stop}, prepínanie módov a vstupy 
pre manuálne PWM hodnoty. Stav \textit{Archív} slúži na zobrazenie historických dát za 
zvolené časové obdobie.

\begin{figure}[H]
    \centering
    \safeimage{obrazky/thingsboard_default.png}{0.95\textwidth}{ThingsBoard dashboard v stave Default}
    \label{fig:thingsboard_default}
\end{figure}

\begin{figure}[H]
    \centering
    \safeimage{obrazky/thingsboard_archiv.png}{0.95\textwidth}{ThingsBoard dashboard v stave Archív}
    \label{fig:thingsboard_archiv}
\end{figure}

\subsection{Obojsmerná komunikácia s ThingsBoardom}

Systém neodosiela do ThingsBoardu iba telemetriu, ale zároveň prijíma aj vybrané riadiace 
príkazy. Použité sú dva mechanizmy. Prvým sú RPC príkazy cez topic 
\path{v1/devices/me/rpc/request/+}, ktoré sa používajú najmä pre tlačidlá 
\textit{Connect}, \textit{Start} a \textit{Stop}. Druhým mechanizmom sú zdieľané atribúty cez 
topic \path{v1/devices/me/attributes}.

Cez zdieľané atribúty sa prenášajú najmä:

\begin{itemize}
    \item \path{control_mode} -- nastavenie regulačného módu,
    \item \path{target_pump_pwm} -- cieľová PWM hodnota pumpy,
    \item \path{target_tec_pwm} -- cieľová PWM hodnota Peltierovho článku.
\end{itemize}

Po prijatí zmeny z ThingsBoardu Flask server aktualizuje svoj interný stav a pri 
najbližšej komunikácii odošle nové hodnoty do ESP32. Zmeny môžu byť zároveň odoslané 
aj do lokálneho webového dashboardu pomocou WebSocket udalosti \path{status_update}.

\subsection{Regulačné módy systému}

Systém podporuje tri regulačné módy, ktoré je možné prepínať z webového rozhrania alebo z 
ThingsBoard dashboardu.

\begin{itemize}
    \item \textbf{Mód 1} -- regulácia pomocou Peltierovho článku pri konštantnom alebo manuálne nastavenom výkone pumpy,
    \item \textbf{Mód 2} -- regulácia pomocou prietoku pumpy pri konštantnom výkone Peltierovho článku,
    \item \textbf{Mód 3} -- hybridná alebo kaskádová regulácia pomocou Peltierovho článku aj pumpy.
\end{itemize}

Zvolený mód sa ukladá do interného stavu servera a následne sa prenáša do ESP32. 
Pre PID reguláciu je v projekte pripravený modul \path{pid_logic.py}, ktorý obsahuje triedu 
\path{PIDController}. Finálne nastavenie regulačných parametrov závisí od výsledkov 
identifikácie a testovania reálnej tepelnej sústavy.
